%%%%%%%%%%%%%%%%%%%%%%%%%%%%%%%%%%%%%%%%%%%%%%%%%%%%%%%%%%%%%%%%%%%%%%
% Lớp 4 | Chủ đề 1. Ôn tập và bổ sung | Số chẵn, số lẻ
% Dạng 4D1?3-1: Dãy số chẵn, số lẻ: đếm số số hạng, số hạng thứ n
%%%%%%%%%%%%%%%%%%%%%%%%%%%%%%%%%%%%%%%%%%%%%%%%%%%%%%%%%%%%%%%%%%%%%%


% nguồn: 2026-11-18_lop4_chu-de-1-tu-luan (nhập 18/11/2026)
\begin{bt}%[4D1N3-1]%[ma:4979b232]
% nguon: Toán tư duy lớp 4 – Bài 3 SGK Toán 4 KNTT
	Có bao nhiêu số chẵn có hai chữ số?
	\loigiai{
		Các số chẵn có hai chữ số là $10, 12, 14, \dots, 98$; hai số liền nhau hơn kém nhau $2$ đơn vị.\\
		Số các số: $(98 - 10) : 2 + 1 = 45$ (số).\\
		Công thức đếm: (số cuối $-$ số đầu) $:$ khoảng cách $+ 1$.
	}
\end{bt}

% nguồn: 2026-11-18_lop4_chu-de-1-tu-luan (nhập 18/11/2026)
\begin{bt}%[4D1H3-1]%[ma:21465ab3]
% nguon: Toán tư duy lớp 4 – Bài 3 SGK Toán 4 KNTT
	Từ $1$ đến $100$ có bao nhiêu số lẻ? Tính tổng các số lẻ đó.
	\loigiai{
		Các số lẻ: $1, 3, 5, \dots, 99$. Số các số lẻ: $(99 - 1) : 2 + 1 = 50$ (số).\\
		Ghép cặp đầu – cuối: $(1 + 99) + (3 + 97) + \dots$, mỗi cặp bằng $100$, có $50 : 2 = 25$ cặp.\\
		Tổng: $100 \times 25 = 2\,500$.
	}
\end{bt}

% nguồn: 2026-11-18_lop4_chu-de-1-tu-luan (nhập 18/11/2026)
\begin{bt}%[4D1V3-1]%[ma:dda35253]
% nguon: Toán tư duy lớp 4 – Bài 3 SGK Toán 4 KNTT
	Bạn An viết liên tiếp các số chẵn từ $2$ đến $200$. Hỏi An đã viết tất cả bao nhiêu chữ số?
	\loigiai{
		Số chẵn có một chữ số: $2, 4, 6, 8$: $4$ số, cần $4$ chữ số.\\
		Số chẵn có hai chữ số: $10$ đến $98$: $(98 - 10) : 2 + 1 = 45$ (số), cần $45 \times 2 = 90$ (chữ số).\\
		Số chẵn có ba chữ số: $100$ đến $200$: $(200 - 100) : 2 + 1 = 51$ (số), cần $51 \times 3 = 153$ (chữ số).\\
		Tất cả: $4 + 90 + 153 = 247$ (chữ số).
	}
\end{bt}

% nguồn: 2026-11-18_lop4_chu-de-1-tu-luan-bo-sung (nhập 18/11/2026)
\begin{bt}%[4D1N3-1]%[ma:53c69b78]
% nguon: Toán tư duy lớp 4 – Bài 3 SGK Toán 4 KNTT
	Tìm số chẵn lớn nhất có bốn chữ số và số lẻ nhỏ nhất có bốn chữ số.
	\loigiai{
		Số lớn nhất có bốn chữ số là $9\,999$ (số lẻ), nên số chẵn lớn nhất có bốn chữ số là $9\,998$.\\
		Số nhỏ nhất có bốn chữ số là $1\,000$ (số chẵn), nên số lẻ nhỏ nhất có bốn chữ số là $1\,001$.
	}
\end{bt}

% nguồn: chuyen_dang (nhập 22/11/2026)
\begin{ex}%[4D1B3-1]%[ma:94ec45cd]
	Cho hai hàm số  $f(x)=\sin2x$ và $g(x)=\cos3x$. Mệnh đề nào sau đây là đúng?
	\choice 
		 {$f(x)$ và $g(x)$ là hai hàm số chẵn}
		{$f(x)$ và  $g(x)$ là hai hàm số lẻ}
		{$f(x)$ là  hàm số chẵn và $g(x)$ là hàm số lẻ}
		{\True $f(x)$ là hàm số lẻ  và $g(x)$ là hàm số chẵn}
	\loigiai{ 
	\begin{itemize}
		\item Xét hàm số $f(x)=\sin  2x$.\\
		Tập  xác định là $\mathscr{D}=\mathbb{R}$.\\
		Với mọi  $x\in\mathscr{D}$ thì $-x\in \mathscr{D}$; $f(-x)=\sin(-2x)=-\sin 2x=-f(x)$.\\
		Vậy hàm số $f(x)=\sin 2x$ là hàm số lẻ. 
		 \item Xét hàm số $f(x)=\cos 3x$.\\
		Tập xác định  là $\mathscr{D}=\mathbb{R}$.\\
		Với mọi $x\in\mathscr{D}$ thì  $-x\in \mathscr{D}$; $f(-x)=\cos(-3x)=\cos 3x=f(x)$.\\
		Vậy hàm số $f(x)=\cos 3x$ là  hàm số chẵn.
	\end{itemize}
	} 
\end{ex}

% nguồn: chuyen_dang (nhập 22/11/2026)
\begin{ex}%[4D1B3-1]%[ma:79b43f5e]
	Trong các  mệnh đề sau, mệnh đề nào là đúng?
	 \choice
		 {Tổng của hai số tự nhiên là một số chẵn khi và chỉ khi cả hai số đều là số chẵn}
		 {Tích của hai số tự nhiên là một số chẵn khi và chỉ khi cả hai số đều là số chẵn}
		{Tổng của hai số tự nhiên là một số  lẻ khi và chỉ khi cả hai số đều là số lẻ}
		{\True Tích của hai số tự nhiên là một số  lẻ khi và chỉ khi cả hai số đều là số lẻ}
	 \loigiai{
		Ta có các nhận  xét sau đây:
		\begin{itemize}
			\item Hai số lẻ vẫn  có tổng là số chẵn.
			\item Chỉ cần ít  nhất một thừa số chẵn thì tích hai số là số chẵn.
			\item Tổng hai số tự nhiên lẻ lại là một số  chẵn.
			\item Tích hai số  tự nhiên lẻ thì trong các thừa số không có thừa số chẵn, hay nói cách khác tất cả các thừa số đều là số lẻ.
		\end{itemize}
	 }
\end{ex}

% nguồn: chuyen_dang (nhập 22/11/2026)
\begin{ex}%[4D1B3-1]%[ma:7c41a82b]
	Cho $a,b$ là hai  số tự nhiên. Mệnh đề nào sau đây \textbf{sai}?
	 \choice
		{Nếu $a,  b$ là các số lẻ thì $ab$ lẻ}
		{\True  Nếu $a$ chẵn và $b$ lẻ thì $ab$ lẻ}
		 {Nếu $a$ và $b$ lẻ thì $a+b$ chẵn}
		{Nếu  $a^2$ lẻ thì $a$ lẻ}
	\loigiai{Mệnh đề “Nếu $a$ chẵn và $b$ lẻ thì $ab$ lẻ”  sai vì tích của một số tự nhiên bất kì với một số tự nhiên chẵn là một số tự nhiên chẵn.}
\end{ex}

% nguồn: chuyen_dang (nhập 22/11/2026)
\begin{ex}%[4D1H3-1]%[ma:149cb405]
	Cho hàm số $f(x)=\sin 2x$ và $g(x)=\tan^2 x$.  Chọn mệnh đề đúng
	 \choice
		{$f(x)$ là  hàm số chẵn, $g(x)$ là hàm số lẻ}
		{\True $f(x)$ là hàm số lẻ,  $g(x)$ là hàm số chẵn}
		{$f(x)$ là hàm số chẵn,  $g(x)$ là hàm số chẵn}
		{$f(x)$  và $g(x)$ đều là hàm số lẻ}
	\loigiai{ 
		Xét  hàm số $f(x)=\sin 2x$.
		\begin{itemize}
			\item Tập xác  định $\mathscr{D}=\mathbb{R}$. Do đó $\forall x\in \mathscr{D}\Rightarrow-x\in \mathscr{D}$.
			\item Ta có  $f\left(-x\right)=\sin \left(-2x\right)=-\sin 2x=-f(x)$.
			 $\Rightarrow f(x)$ là hàm số lẻ.
		\end{itemize}
		Xét hàm số $g(x)=\tan^2  x$.
		\begin{itemize}
			\item Tập xác định $\mathscr{D}=\mathbb{R}\setminus \left\{\dfrac{\pi}{2}+k\pi, \;k\in \mathbb{Z}\right\}$. Do đó $\forall x\in \mathscr{D}\Rightarrow-x\in \mathscr{D}$. 
			\item Ta  có $g\left(-x\right)=\left[\tan \left(-x\right)\right]^2=\left(-\tan x\right)^2=\tan^2 x=g(x)$.
			$\Rightarrow g(x)$ là hàm số  chẵn.
		\end{itemize}
	} 
\end{ex}

% nguồn: chuyen_dang (nhập 22/11/2026)
\begin{ex}%[4D1H3-1]%[ma:824b26bd]
	Cho hàm số  $f(x)=\sin 2x$ và $g(x)=\tan ^2x$. Chọn mệnh đề đúng?
	\choice 
		 {$f(x)$ là hàm số chẵn, $g(x)$ là hàm số lẻ}
		{\True $f(x)$ là hàm số lẻ,  $g(x)$ là hàm số chẵn}
		{$f(x)$ là hàm số chẵn, $g(x)$ là hàm số  chẵn}
		{$f(x)$ và $g(x)$  đều là hàm số lẻ}
	 \loigiai{
		Hàm  số $f(x)$ có tập xác định là $\mathscr{D}=\mathbb{R}\Rightarrow \forall x\in \mathscr{D}$ thì $-x\in \mathscr{D}$.\\
		$\forall x\in \mathscr{D}\colon f (-x)=\sin  \left(-2x\right)=-\sin 2x=-f(x)$ $\Rightarrow$ hàm số $f(x)$ là hàm số lẻ.\\
		Hàm số $g(x)$ có tập xác định là $\mathscr{D}=\mathbb{R}\setminus \left\{\dfrac{\pi}{2}+k\pi, k\in {\bf  \mathbb{Z}}\right\}$ $\Rightarrow \forall x\in \mathscr{D}$ thì $-x\in \mathscr{D}$.\\
		$\forall x\in \mathscr{D}\colon g (-x)=\tan ^2 (-x)=\tan  ^2x=g(x)$ $\Rightarrow$ hàm số $g(x)$ là hàm số chẵn.
	} 
\end{ex}

% nguồn: chuyen_dang (nhập 22/11/2026)
\begin{ex}%[4D1H3-1]%[ma:1b5eccd6]
	Cho hai hàm số $f(x)=\sin 2x$ và $g(x)=\cos 3x$. Chọn mệnh đề  đúng
	 \choice
		{$f$ là hàm  số chẵn và $g$ là hàm số lẻ}
		 {$f$ và $g$ là hai hàm số chẵn}
		{$f$ và  $g$ là hai hàm số lẻ}
		{\True $f$ là hàm số  lẻ và $g$ là hàm số chẵn}
	 \loigiai{
		Tập xác định của hai hàm số  là $\mathscr{D}=\mathbb{R}$.\\
		 Ta có $f (-x)=\sin (-2x)=-\sin 2x=-f(x)\Rightarrow$ $f$ là hàm số lẻ.\\
		$g (-x)=\cos (-3x)=\cos 3x=g(x)$ $\Rightarrow$ $g$ là hàm  số chẵn.
	 }
\end{ex}

% nguồn: chuyen_dang (nhập 22/11/2026)
\begin{ex}%[4D1B3-1]%[ma:89750980]
	Cho $x$ là số tự nhiên. Phủ định của mệnh đề  “$\forall x$ chẵn, $x^2+x$ là số chẵn” là mệnh đề:
	\choice 
		{$\exists x$ lẻ,  $x^2+x$ là số lẻ}
		{$\exists x$  lẻ, $x^2+x$ là số chẵn}
		{$\forall x$ lẻ, $x^2+x$  là số lẻ}
		{\True  $\exists x$ chẵn, $x^2+x$ là số lẻ}
	\loigiai{ 
	Mệnh đề phủ định là  “$\exists x$ chẵn, $x^2+x$ là số lẻ”.}
\end{ex}

% nguồn: chuyen_dang (nhập 22/11/2026)
\begin{bt}%[4D1K3-1]%[ma:8bfb6725]
	Chứng minh bằng  phản chứng. Cho $n \in \mathbb{N}$, nếu $5n+5$ là số lẻ thì $n$ là số chẵn.
	 \loigiai{
		Giả sử $n$ là số lẻ. Suy ra $n=2k+1,~\forall  k \in \mathbb{N}$.\\
		Khi đó $5n+5=5(2k+1)+5=10k+10=10(k+1)$ (trái giả thiết $5n+5$ là  số lẻ).\\
		Vậy phải có  $n$ là số chẵn.
	 }
\end{bt}
